\documentclass[12pt,a4paper]{article}

\usepackage[french]{babel}
\usepackage[utf8]{inputenc}
\usepackage[T1]{fontenc}
\usepackage{fouriernc}
\usepackage{geometry}
\usepackage{array}
\usepackage{booktabs}
\usepackage{tabularx}
\usepackage[table]{xcolor}
\usepackage{amsmath}
\usepackage{fancyhdr}
\usepackage[most]{tcolorbox}
\usepackage{enumitem}

\geometry{top=2.2cm, bottom=2cm, left=2cm, right=2cm, headheight=15pt}

%------------------------------------------------------------
% Couleurs
%------------------------------------------------------------
\definecolor{bluedark}{RGB}{50, 100, 160}
\definecolor{bluemid}{RGB}{170, 200, 235}
\definecolor{grisfond}{RGB}{245, 245, 245}
\definecolor{vertcorr}{RGB}{34, 120, 60}
\definecolor{vertlight}{RGB}{220, 245, 225}
\definecolor{vertmid}{RGB}{150, 210, 170}

%------------------------------------------------------------
% En-tête / pied de page
%------------------------------------------------------------
\pagestyle{fancy}
\fancyhf{}
\lhead{\small Physique-Chimie -- Première Bac Pro Microtechniques}
\rhead{\small \textbf{CORRIGÉ} -- Les solutions aqueuses}
\cfoot{\thepage}
\renewcommand{\headrulewidth}{0.5pt}

%------------------------------------------------------------
% Boîtes tcolorbox
%------------------------------------------------------------
\tcbset{
  stylepartie/.style={
    colback=bluedark, colframe=bluedark, coltext=white,
    boxrule=0pt, arc=3pt,
    left=6pt, right=6pt, top=4pt, bottom=4pt,
    before skip=10pt, after skip=6pt,
  },
  stylecorr/.style={
    colback=vertlight, colframe=vertmid,
    boxrule=0.8pt, arc=4pt,
    left=8pt, right=8pt, top=6pt, bottom=6pt,
    before skip=4pt, after skip=6pt,
  },
}

%------------------------------------------------------------
% Commandes
%------------------------------------------------------------
\newcounter{numquestion}
\newcommand{\question}[1]{%
  \stepcounter{numquestion}%
  \vspace{8pt}%
  \noindent{\color{bluedark}$\blacktriangleright$~\textbf{Question~\thenumquestion.}}~#1%
  \vspace{2pt}%
}

\newcommand{\repcorr}[1]{%
  \begin{tcolorbox}[stylecorr]
    {\color{vertcorr} #1}
  \end{tcolorbox}%
}

%------------------------------------------------------------
\begin{document}
%------------------------------------------------------------

%============================================================
% TITRE
%============================================================
\begin{tcolorbox}[
  colback=white, colframe=bluedark, boxrule=1.5pt, arc=4pt,
  before skip=0pt, after skip=8pt,
  left=10pt, right=10pt, top=8pt, bottom=8pt,
]
  \begin{center}
    {\color{bluedark}\Large\bfseries Dissolution et Dilution}\\[4pt]
    {\large Activité documentaire -- \textbf{\color{vertcorr}CORRIGÉ}}\\[6pt]
    \hrule height 0.4pt \vspace{6pt}
    \textbf{Durée :} 1\,h\,00 \hfill \textbf{Niveau :} Première Bac Pro \hfill \textbf{Total : 20 points}
  \end{center}
\end{tcolorbox}

\vspace{6pt}

%============================================================
% PARTIE A
%============================================================
\begin{tcolorbox}[stylepartie]
  Partie A -- La dissolution \hfill {\normalfont\small (20 min -- 5 pts)}
\end{tcolorbox}

%------------------------------------------------------------
\question{Identifier dans le texte le \textbf{soluté} et le \textbf{solvant} de la solution préparée.}

\repcorr{
  \textbf{Soluté :} chlorure de sodium (NaCl) — c'est le solide dissous.\\[3pt]
  \textbf{Solvant :} eau distillée — c'est le liquide dans lequel on dissout.
}

%------------------------------------------------------------
\question{Calculer la masse molaire $M(\text{NaCl})$ du chlorure de sodium (en g/mol).}

\repcorr{
  On additionne les masses molaires atomiques de chaque atome :
  \[
    M(\text{NaCl}) = M(\text{Na}) + M(\text{Cl}) = 23 + 35{,}5 = \boxed{58{,}5~\text{g/mol}}
  \]
}

%------------------------------------------------------------
\question{Calculer la concentration massique $C_m$ (en g/L) de la solution.}

\repcorr{
  \textbf{Conversion du volume :} $V = 500~\text{mL} = 0{,}500~\text{L}$
  \[
    C_m = \dfrac{m}{V} = \dfrac{4{,}5~\text{g}}{0{,}500~\text{L}} = \boxed{9~\text{g/L}}
  \]
}

\vspace{4pt}

%============================================================
% PARTIE B
%============================================================
\begin{tcolorbox}[stylepartie]
  Partie B -- La dilution \hfill {\normalfont\small (20 min -- 8 pts)}
\end{tcolorbox}

%------------------------------------------------------------
\question{Expliquer, en une phrase, ce qu'est une \textbf{dilution} et ce qui est conservé lors de cette opération.}

\repcorr{
  La dilution consiste à préparer une solution moins concentrée (\textbf{solution fille}) à
  partir d'une solution plus concentrée (\textbf{solution mère}) en ajoutant du solvant (eau
  distillée). La \textbf{quantité de soluté} (et donc la masse de soluté) est conservée.
}

%------------------------------------------------------------
\question{Calculer le \textbf{facteur de dilution} $f$ entre la solution mère et la solution fille.}

\repcorr{
  \[
    f = \dfrac{C_1}{C_2} = \dfrac{80~\text{g/L}}{10~\text{g/L}} = \boxed{8}
  \]
  La solution mère est 8 fois plus concentrée que la solution fille.
}

%------------------------------------------------------------
\question{En utilisant la relation $C_1 V_1 = C_2 V_2$, calculer le volume $V_1$ (en mL) de solution mère à prélever.}

\repcorr{
  On isole $V_1$ :
  \[
    V_1 = \dfrac{C_2 \times V_2}{C_1} = \dfrac{10 \times 2{,}0}{80} = 0{,}25~\text{L} = \boxed{250~\text{mL}}
  \]
  \textbf{Vérification :} $C_1 V_1 = 80 \times 0{,}25 = 20$ \quad ; \quad $C_2 V_2 = 10 \times 2{,}0 = 20$ {\color{vertcorr}$\surd$}
}

%------------------------------------------------------------
\question{En vous aidant du document~3, nommer le matériel que le technicien doit utiliser pour prélever $V_1$ avec précision, puis pour préparer la solution fille. Justifier.}

\repcorr{
  \begin{itemize}[noitemsep]
    \item Pour prélever $V_1 = 250$~mL avec précision : \textbf{pipette graduée} — elle permet de
    prélever un volume quelconque avec une bonne précision.
    \item Pour préparer la solution fille de volume exact $V_2 = 2{,}0$~L : \textbf{fiole jaugée de
    2,0~L} — elle garantit un volume exact grâce au trait de jauge gravé.
  \end{itemize}
}

\vspace{4pt}

%============================================================
% PARTIE C
%============================================================
\begin{tcolorbox}[stylepartie]
  Partie C -- Concentration molaire et bilan de synthèse \hfill {\normalfont\small (20 min -- 7 pts)}
\end{tcolorbox}

%------------------------------------------------------------
\question{Calculer la \textbf{concentration molaire} $C$ (en mol/L) de la solution d'acide citrique du bain.}

\repcorr{
  \[
    C = \dfrac{C_m}{M} = \dfrac{96~\text{g/L}}{192~\text{g/mol}} = \boxed{0{,}5~\text{mol/L}}
  \]
}

%------------------------------------------------------------
\question{Calculer la \textbf{quantité de matière} $n$ d'acide citrique contenue dans le bain de 5,0~L.}

\repcorr{
  \[
    n = C \times V = 0{,}5~\text{mol/L} \times 5{,}0~\text{L} = \boxed{2{,}5~\text{mol}}
  \]
}

%------------------------------------------------------------
\question{\textbf{Décision :} Le bain doit être renouvelé si $C < 0{,}25$~mol/L. Conclure si le bain est encore utilisable.}

\repcorr{
  $C = 0{,}5~\text{mol/L} > 0{,}25~\text{mol/L}$ (seuil de remplacement).\\[3pt]
  $\Rightarrow$ Le bain est \textbf{encore utilisable}, il n'a pas besoin d'être renouvelé.
}

%------------------------------------------------------------
\question{\textbf{Synthèse :} Expliquer la \textbf{différence fondamentale} entre une dissolution et une dilution. Donner un exemple de chaque opération tiré de l'activité.}

\repcorr{
  \begin{itemize}[noitemsep]
    \item \textbf{Dissolution :} on part d'un \textbf{soluté solide} que l'on introduit dans un
    solvant pour obtenir une solution. \textit{Exemple dans l'activité :} préparation de la solution
    saline en dissolvant 4,5~g de NaCl dans de l'eau (Doc~1).
    \item \textbf{Dilution :} on part d'une \textbf{solution déjà préparée} (solution mère) à
    laquelle on ajoute du solvant pour diminuer la concentration. \textit{Exemple dans l'activité :}
    dilution du dégraissant industriel ($C_1 = 80$~g/L) pour obtenir $C_2 = 10$~g/L (Doc~2).
  \end{itemize}
  Dans les deux cas on ajoute de l'eau, mais le point de départ est différent : \textbf{solide} pour
  la dissolution, \textbf{solution concentrée} pour la dilution.
}

\vspace{10pt}
\begin{center}
  \color{bluedark}\rule{0.5\linewidth}{0.6pt}
\end{center}

\end{document}
